\documentclass[12pt]{article}
\usepackage[utf8]{inputenc}
\usepackage[T1]{fontenc}
\usepackage{hyperref}
\title{Moo Language Resources for K-12 Curriculum Development}
\author{Compiled by Claude Sonnet 5.5 (claude-sonnet-5-5) for Andr\'as Kornai}
\date{}
\begin{document}
\maketitle

\section{Introduction}
Moo (Gomu, Gwomu, ISO 639-3 gwg) is a Bikwin (Adamawa) language of Taraba State, Nigeria. Population estimates differ (about 5,000 in an older figure, 9,100 in another listing). It should not be confused with Gwa (gwb). The search results offered only catalogue entries; no educational or reference materials were found.

\section{Primary Educational Resources}
None found.

\section{Linguistic and Reference Materials}
\begin{itemize}
\item Wikipedia article on Moo: \url{https://en.wikipedia.org/wiki/Moo_language}
\item ASJP language page (word list data; contents not examined): \url{https://asjp.clld.org/languages/MOO}
\end{itemize}

\section{Religious and Community Resources}
\begin{itemize}
\item Joshua Project listing ``Gwomu in Nigeria'' (listing seen in search results; contents not examined, and its match to Moo is presumed from the name): \url{https://joshuaproject.net/people_groups/12021}
\end{itemize}

\section{Digital and Audio Resources}
\begin{itemize}
\item Global Recordings Network page for gwg, which states that no recordings are currently available in this language: \url{https://globalrecordings.net/en/language/gwg}
\end{itemize}

\section{Curriculum Development Context}
There is no evidence of orthography, literacy materials or schooling in Moo. Curriculum work would begin from documentation and community consultation.

\end{document}